В Кыргызстане клиент дизайн-студии обычно утверждает проект по 360\textdegree-панорамам и рендерам. Ракурс выбирает дизайнер, а широкоугольная камера и <<идеальный>> свет делают комнаты визуально больше и светлее, чем в реальности. В результате клиент неверно оценивает размеры, освещение и то, как открываются двери и мебель, а ошибки обнаруживаются только во время ремонта.

В работе представлен JK3DStudio VR "--- инструмент, который превращает готовый дизайн-проект из 3ds Max в интерактивную VR-презентацию. Клиент ходит по будущей квартире в масштабе 1:1, открывает двери рукой, включает свет по комнатам, меняет время суток и цветовую температуру. ИИ-гид на русском и английском рассказывает о проекте и говорит только факты из паспорта проекта, утверждённого дизайнером.

Система состоит из трёх частей: плагина для редактора Unreal Engine 5.8, который полуавтоматически готовит сцену по стандарту имён, VR-приложения для очков PICO 4 Ultra и серверной части AI Gateway (FastAPI, PostgreSQL) для генерации речи, проверки фактов и озвучки. За рубежом подобные интерактивные туры делают вручную, их стоимость "--- от 10\,000 долларов США; аналогичной услуги для дизайн-проектов в Кыргызстане не найдено. Новизна "--- быстрая подготовка, честные масштаб и свет и ИИ-гид, опирающийся на факты, с адаптацией к местному рынку. На первой предзащите демонстрируются рабочий прототип VR-приложения и протестированная серверная часть.

\medskip\noindent\textbf{Ключевые слова:} виртуальная реальность, дизайн интерьера, Unreal Engine, ИИ-гид, большие языковые модели, Кыргызстан.
